\section{Package Manager}

\begin{frame}[fragile]{Introduction to package manager}
	\begin{itemize}
		\item A package manager or package management system (PMS) is a collection of software tools that automates the process of installing, upgrading, configuring, and removing computer programs for a computer in a consistent manner.
	\end{itemize}
\end{frame}

\begin{frame}[fragile]{Problem solved}
	\begin{itemize}
		\item \textbf{Dependency Hell}: Different software packages require different, and sometimes conflicting, versions of the same shared libraries. Package managers solve this by managing and allowing for the coexistence of multiple library versions.
		\item \textbf{Manual Installation}: Package managers eliminate the need for users to manually download, compile, and install software, which can be a complex and time-consuming process.
		\item \textbf{Synchronization Issues}: They ensure that the list of installed software is always consistent and up-to-date with a central database, preventing conflicts and missing prerequisites that could arise from manual interventions.
	\end{itemize}
\end{frame}

\begin{frame}[fragile]{How It Works}

	A package manager operates by interacting with three key components:

	\begin{itemize}
		\item \textbf{Packages}: Packages are the fundamental units of a PMS. A package is a file that contains the application, its necessary files, and metadata like the name, version, and dependencies.
		\item \textbf{Repositories}: These are centralized locations or servers where packages are stored. A package manager downloads packages from these repositories.
		\item \textbf{Local Database}: The package manager maintains a local database on the user's system. This database keeps a record of all installed packages, their versions, and their dependencies.
	\end{itemize}
\end{frame}

\begin{frame}[fragile]{Mirror}
	\begin{itemize}
		\item The mirror/source config of package defines where your package manager fetch remote packages
		\item Can be customized to improve download speed and availability
		\item Many good quality sources include: Tsinghua TUNA, USTC, SJTUG...
	\end{itemize}

	This link from CERNET would give you configuration steps on pretty much everything they provide: \href{https://help.mirrors.cernet.edu.cn/}{CERNET mirrors guide}
\end{frame}

\begin{frame}{Try it out!}

	In the WSL Install Guide, we provide commands for you to install VS Code through package managers.

	Follow the \textbf{VS Code Setup} section in \texttt{wsl.pdf} and try to use the commands!

	Note: For Ubuntu users, \textbf{snap} is \textbf{not recommended}

\end{frame}